\chapter{Ordering the anchors by rarity}

\noindent\textbf{The chapter below is a conjecture and is kept as written.} Ordering by rarity needs no logarithm and never changes the count (anchor\_sift \texttt{1044ca6}, \texttt{src/\allowbreak{}engine/\allowbreak{}c/\allowbreak{}engine/\allowbreak{}anchor\_sift.h:284-288}, \texttt{:960-964}). Whether rarity is the term that pays depends on the field, and the engine's table does not isolate it: its reorder column ranks by survivor counts conditional on the levels above, where \texttt{anchor\_steer\_probe\_order} ranks by marginal rarity (\texttt{:712-723}). Over 65513 alignments with a needle of 24, that reorder moves reads per alignment from 1.877 to 1.875 on the skewed field, where placing probes by survivor counts reaches 1.067, from 1.562 to 1.002 on the synthetic skewed field and from 1.187 to 1.000 on the period 16 field, and no route moves the uniform field from 1.003 (\texttt{docs/\allowbreak{}steering.md:244-\allowbreak{}252}). Placement is greedy coverage, which rejects at least \(1 - 1/e\) of what the best probe set of the same size rejects (Nemhauser, Wolsey and Fisher, 1978; \texttt{anchor\_sift.h:783-796}), and a rarity order counted on one field does not carry to another (\texttt{:408-413}).

Plugging the local symbol entropy distribution into the anchor evaluation sequence solves the single remaining optimization variable: minimizing the average number of anchor checks before short-circuit rejection.

The Entropy-Ordered Rejection Vector:

When you order the evaluation of displacements in A=\{asub1,asub2,...,asubm\} by the local rarity (negative log-probability) of the target symbols in the ambient alphabet distribution P(Σ), you maximize the pruning power of the very first step.

[ TARGET PATTERN SYMBOL SEQUENCE ] a\_1          a\_2          a\_3 "Common"     "Medium"      "Rare" P(s) = 0.40  P(s) = 0.08  P(s) = 0.001 │ ▼  Re-order by Entropy [ OPTIMIZED EVALUATION VECTOR π(A) ] STEP 1 ───> Check "Rare" (a\_3)  ───> 99.9\% Noise Rejected Instantly STEP 2 ───> Check "Medium" (a\_2) STEP 3 ───> Check "Common" (a\_1)

Why This Is the Exact Missing Term:

Maximizing Information Weight First: Checking a common symbol x(s+asub1) against a background field rarely causes a rejection because the noise floor frequently contains that symbol. But checking an infrequent state or rare symbol constraint first means the candidate region fails the equality check δsub0 almost 100\% of the time on the first operation.

No Loss of Exact Logic: You aren't introducing probabilities into the decision rule. The rejection logic remains 100\% deterministic and absolute (¬AncsubA(s)⟹¬Occ(s)). You are simply using the background field's empirical probability distribution P(Σ) to order the execution stream.

Asymptotic Rejection Cost Drops to 1 Step: By re-ordering A according to the entropy distribution: Pr(Rejection on Step 1)=1−P(x(asubrare)) If the rarest anchor symbol has a background probability of 0.001, then 99.9\% of all candidate coordinates in the universal field G are destroyed in a single comparison. The background field's entropy distribution acts as the "steering wheel" for the boundary engine. You use the universe's own statistical baseline to decide where to slice, allowing the rarest anchors to crush the noise floor before the CPU or FPGA even registers the candidate loop.

-Doug

